\chapter{Zusammenfassung}

In diesem Dokument finden Sie eine Sammlung von wertvollen Tipps und Empfehlungen, die Ihnen beim Verfassen Ihrer schriftlichen Abschlussarbeit am Lehrstuhl für Medizinischen Informationstechnik (MEDIT) helfen sollen. Es geht dabei um verschiedene wichtige Aspekte, die Sie bei der Erstellung Ihrer Arbeit berücksichtigen sollten, um eine qualitativ hochwertige wissenschaftliche Arbeit zu gewährleisten.

Gleichzeitig ist dieses Dokument auch eine Vorlage, die direkt für die schriftliche Arbeit genutzt werden kann.

Das Dokument beginnt mit Hinweisen zum Ablauf der Arbeit, um Ihnen einen Überblick über die einzelnen Phasen der Erstellung einer Abschlussarbeit zu geben. Anschließend wird der Aufbau der Arbeit behandelt, wobei auf die verschiedenen Abschnitte wie Einleitung, Hauptteil und Schluss eingegangen wird. Auch die inhaltliche Gestaltung der Arbeit wird thematisiert, wobei es darum geht, wie Sie Ihre Forschungsfrage präzise formulieren, Ihre Methodik beschreiben und Ihre Ergebnisse darstellen und interpretieren. Nicht zuletzt geben wir einige Hinweise zu häufig verwendeten Tools.

Ein weiterer wichtiger Aspekt, der in dem Dokument behandelt wird, ist der Umgang mit Quellen. Hierbei geht es um die korrekte Zitierung und Anerkennung von fremden Gedanken und Ideen, um Plagiate zu vermeiden. Auch die Bedeutung einer sorgfältigen Recherche und Auswahl von Quellen wird hervorgehoben.

Insgesamt dient das Dokument als umfassender Leitfaden für Studierende und gibt wertvolle Tipps und Empfehlungen, um den Schreibprozess zu erleichtern und die Qualität der Arbeit zu verbessern.



\cleardoubleemptypage
\thispagestyle{empty}